% ============================================================
\chapter{DỮ LIỆU VÀ PHƯƠNG PHÁP THỰC NGHIỆM}
\label{ch:du_lieu}
% ============================================================

\section{Bộ dữ liệu}

Báo cáo sử dụng bộ dữ liệu \textbf{Forest Inspection} (mô phỏng bằng AirSim), công bố trên Zenodo~\cite{forestdataset}. Dữ liệu gồm \textbf{18 chuỗi ảnh} (sequence), kết hợp cả hai điều kiện thời tiết \textit{sunny} (9 chuỗi từ seq1--9) và \textit{overcast} (9 chuỗi từ seq10--18), được thu ở các tổ hợp góc chụp (pitch $\in\{0^\circ, -60^\circ, -90^\circ\}$) và độ cao bay ($\{30, 50, 80\}$ m). Tổng cộng bộ dữ liệu gồm **26.242** cặp ảnh (RGB + nhãn phân đoạn) ở độ phân giải gốc $1280\times 768$. Mỗi điểm ảnh được gán một trong \textbf{11 lớp}: Sky, Deciduous\_trees, Coniferous\_trees, Fallen\_trees, Dirt\_ground, Ground\_vegetation, Rocks, Building, Fence, Car và Empty (nền/không xác định).

\begin{table}[H]
\centering
\caption{Đặc tả 9 chuỗi ảnh theo góc chụp và độ cao.}
\label{tab:sequences}
\small
\begin{tabular}{c c c c c l}
\toprule
\textbf{Seq} & \textbf{Pitch} & \textbf{Độ cao (m)} & \textbf{Số ảnh} & \textbf{Đặc điểm} \\
\midrule
seq1 & $0^\circ$   & 30 & 1.468 & Ngang, thấp \\
seq2 & $-60^\circ$ & 30 & 1.501 & Chéo, thấp \\
seq3 & $-90^\circ$ & 30 & 1.469 & Thẳng xuống, thấp \\
seq4 & $0^\circ$   & 50 & 1.432 & Ngang, trung bình \\
seq5 & $-60^\circ$ & 50 & 1.447 & Chéo, trung bình \\
seq6 & $-90^\circ$ & 50 & 1.437 & Thẳng xuống, trung bình \\
seq7 & $0^\circ$   & 80 & 1.436 & Ngang, cao \\
seq8 & $-60^\circ$ & 80 & 1.464 & Chéo, cao \\
seq9 & $-90^\circ$ & 80 & 1.397 & Thẳng xuống, cao \\
\bottomrule
\end{tabular}
\end{table}

Quá trình kiểm định bảng màu (\textit{palette}) cho thấy $100\%$ điểm ảnh ánh xạ
chính xác về 11 lớp trên cả 9 chuỗi --- không có điểm ảnh ngoài bảng màu, khẳng định
chất lượng nhãn.

\section{Phân tích mất cân bằng lớp}
\label{sec:imbalance_data}

Bảng~\ref{tab:imbalance} thống kê tần suất điểm ảnh theo lớp (trên tập huấn luyện) và
trọng số nghịch tần suất tương ứng. Chênh lệch giữa lớp lớn nhất
(Ground\_vegetation) và lớp nhỏ nhất (Fence) lên tới \textbf{2.761 lần} --- mức mất
cân bằng rất nghiêm trọng.

\begin{table}[H]
\centering
\caption{Phân bố tần suất lớp và mức mất cân bằng (tập huấn luyện).}
\label{tab:imbalance}
\small
\begin{tabular}{c l r r r}
\toprule
\textbf{Hạng} & \textbf{Lớp} & \textbf{Tần suất (\%)} & \textbf{Tỷ lệ so lớp min} & \textbf{Trọng số (nghịch tần suất)} \\
\midrule
1  & Ground\_vegetation & 32{,}40 & 2.761$\times$ & 0{,}003 \\
2  & Coniferous\_trees  & 27{,}85 & 2.373$\times$ & 0{,}003 \\
3  & Sky                & 20{,}33 & 1.733$\times$ & 0{,}004 \\
4  & Deciduous\_trees   & 15{,}11 & 1.288$\times$ & 0{,}006 \\
5  & Dirt\_ground       & 1{,}83  & 156$\times$   & 0{,}046 \\
6  & Rocks              & 1{,}74  & 148$\times$   & 0{,}049 \\
7  & Fallen\_trees      & 0{,}52  & 44$\times$    & 0{,}164 \\
8  & Building           & 0{,}18  & 15$\times$    & 0{,}478 \\
9  & Car                & 0{,}04  & 3{,}6$\times$ & 2{,}026 \\
10 & \textbf{Fence}     & \textbf{0{,}01} & \textbf{1$\times$} & \textbf{7{,}221} \\
\bottomrule
\end{tabular}
\end{table}

\noindent\textbf{Nhận định:} bốn lớp \textit{Fallen\_trees, Building, Car, Fence}
đều dưới $1\%$; riêng Fence ($0{,}01\%$) và Car ($0{,}04\%$) là \textit{cực hiếm}.
Đây chính là các lớp được kỳ vọng hưởng lợi nhiều nhất từ chiến lược xử lý mất cân
bằng.

\section{Sự thay đổi miền (domain gap)}

Phân bố lớp biến đổi mạnh theo góc chụp. Đáng chú ý, lớp \textbf{Sky chỉ xuất hiện ở
các chuỗi $pitch = 0^\circ$} (chiếm $36$--$58\%$) và \textbf{biến mất hoàn toàn}
($0\%$) ở các chuỗi $-60^\circ$ và $-90^\circ$. Vì vậy, IoU của Sky phụ thuộc hoàn
toàn vào việc tập kiểm thử có chứa chuỗi chụp ngang hay không --- một yếu tố cần lưu
ý khi diễn giải kết quả (Chương~\ref{ch:thao_luan}).

\begin{figure}[H]
\centering
\resizebox{0.95\textwidth}{!}{%
\begin{tikzpicture}[xscale=1.1, yscale=1.0]
  % Vẽ hệ trục tọa độ
  \draw[->, thick] (0,0) -- (11,0) node[right] {\textbf{Lớp đối tượng}};
  \draw[->, thick] (0,0) -- (0,5.5) node[above] {\textbf{Tần suất (\% - thang đo log)}};
  
  % Vẽ các vạch chia thang log (0.01%, 0.1%, 1%, 10%, 100%)
  \draw[dashed, gray!40] (0,1) -- (10.5,1) node[right] {\tiny 0.01\%};
  \draw[dashed, gray!40] (0,2) -- (10.5,2) node[right] {\tiny 0.1\%};
  \draw[dashed, gray!40] (0,3) -- (10.5,3) node[right] {\tiny 1\%};
  \draw[dashed, gray!40] (0,4) -- (10.5,4) node[right] {\tiny 10\%};
  \draw[dashed, gray!40] (0,5) -- (10.5,5) node[right] {\tiny 100\%};
  
  % Nhãn trục Y
  \node[left] at (0,1) {\scriptsize 0.01\%};
  \node[left] at (0,2) {\scriptsize 0.1\%};
  \node[left] at (0,3) {\scriptsize 1.0\%};
  \node[left] at (0,4) {\scriptsize 10\%};
  \node[left] at (0,5) {\scriptsize 100\%};
  
  % Vẽ các cột dữ liệu (mapped log-scale manually: y = log10(freq) + 3)
  % Sky (15.21%) -> y = 4.18
  \draw[fill=blue!40, draw=black] (0.7,0) rectangle (1.3, 4.18);
  \node[above, font=\tiny\bfseries] at (1, 4.18) {15.21\%};
  \node[below, rotate=45, anchor=north east, font=\scriptsize] at (1,0) {Sky};
  
  % Deciduous (16.36%) -> y = 4.21
  \draw[fill=blue!40, draw=black] (1.7,0) rectangle (2.3, 4.21);
  \node[above, font=\tiny\bfseries] at (2, 4.21) {16.36\%};
  \node[below, rotate=45, anchor=north east, font=\scriptsize] at (2,0) {Deciduous};
  
  % Coniferous (28.33%) -> y = 4.45
  \draw[fill=blue!40, draw=black] (2.7,0) rectangle (3.3, 4.45);
  \node[above, font=\tiny\bfseries] at (3, 4.45) {28.33\%};
  \node[below, rotate=45, anchor=north east, font=\scriptsize] at (3,0) {Coniferous};
  
  % Fallen trees (0.57%) -> y = 2.76
  \draw[fill=red!40, draw=black] (3.7,0) rectangle (4.3, 2.76);
  \node[above, font=\tiny\bfseries] at (4, 2.76) {0.57\%};
  \node[below, rotate=45, anchor=north east, font=\scriptsize] at (4,0) {Fallen\_trees};
  
  % Dirt ground (1.99%) -> y = 3.30
  \draw[fill=blue!40, draw=black] (4.7,0) rectangle (5.3, 3.30);
  \node[above, font=\tiny\bfseries] at (5, 3.30) {1.99\%};
  \node[below, rotate=45, anchor=north east, font=\scriptsize] at (5,0) {Dirt\_ground};
  
  % Ground vegetation (35.41%) -> y = 4.55
  \draw[fill=blue!40, draw=black] (5.7,0) rectangle (6.3, 4.55);
  \node[above, font=\tiny\bfseries] at (6, 4.55) {35.41\%};
  \node[below, rotate=45, anchor=north east, font=\scriptsize] at (6,0) {Ground\_veg};
  
  % Rocks (1.88%) -> y = 3.27
  \draw[fill=blue!40, draw=black] (6.7,0) rectangle (7.3, 3.27);
  \node[above, font=\tiny\bfseries] at (7, 3.27) {1.88\%};
  \node[below, rotate=45, anchor=north east, font=\scriptsize] at (7,0) {Rocks};
  
  % Building (0.19%) -> y = 2.28
  \draw[fill=red!40, draw=black] (7.7,0) rectangle (8.3, 2.28);
  \node[above, font=\tiny\bfseries] at (8, 2.28) {0.19\%};
  \node[below, rotate=45, anchor=north east, font=\scriptsize] at (8,0) {Building};
  
  % Fence (0.01%) -> y = 1.00
  \draw[fill=red!40, draw=black] (8.7,0) rectangle (9.3, 1.00);
  \node[above, font=\tiny\bfseries] at (9, 1.00) {0.01\%};
  \node[below, rotate=45, anchor=north east, font=\scriptsize] at (9,0) {Fence};
  
  % Car (0.05%) -> y = 1.70
  \draw[fill=red!40, draw=black] (9.7,0) rectangle (10.3, 1.70);
  \node[above, font=\tiny\bfseries] at (10, 1.70) {0.05\%};
  \node[below, rotate=45, anchor=north east, font=\scriptsize] at (10,0) {Car};
  
  % Chú giải
  \draw[fill=white, draw=black] (7.2,3.6) rectangle (10.5, 4.8);
  \draw[fill=blue!40, draw=black] (7.4,4.4) rectangle (7.8, 4.6);
  \node[right, font=\tiny] at (7.9, 4.5) {Lớp đa số ($\ge 1\%$)};
  \draw[fill=red!40, draw=black] (7.4,3.8) rectangle (7.8, 4.0);
  \node[right, font=\tiny] at (7.9, 3.9) {Lớp thiểu số ($< 1\%$)};
\end{tikzpicture}}
\caption{Phân bố tần suất điểm ảnh của 10 lớp trên toàn bộ tập dữ liệu huấn luyện kết hợp (sunny+overcast), biểu diễn theo thang đo logarithmic nhằm làm rõ sự chênh lệch lớn ở các lớp cực hiếm.}
\label{fig:class_distribution}
\end{figure}

\section{Chia tập dữ liệu}
\label{sec:split}

Toàn bộ các thí nghiệm trong báo cáo dùng cách chia theo chuỗi như sau:

\begin{table}[H]
\centering
\caption{Cách chia tập dữ liệu kết hợp (sunny+overcast) dùng trong các thí nghiệm của báo cáo.}
\label{tab:split}
\small
\begin{tabular}{l l r}
\toprule
\textbf{Tập} & \textbf{Các chuỗi ảnh thành phần} & \textbf{Số lượng ảnh} \\
\midrule
Huấn luyện (Train) & seq1, seq3, seq4, seq5, seq8, seq9 (sunny) & 17.454 \\
                   & seq10, seq12, seq14, seq15, seq16, seq17 (overcast) & \\
\midrule
Kiểm định (Val)    & seq2 (sunny), seq13, seq18 (overcast) & 4.297 \\
\midrule
Kiểm thử (Test)    & seq6, seq7 (sunny), seq11 (overcast) & 4.491 \\
\midrule
\multicolumn{2}{l}{\textbf{Tổng cộng toàn bộ}} & \textbf{26.242} \\
\bottomrule
\end{tabular}
\end{table}

Đây là cách chia \textbf{theo chuỗi} (cross-sequence) chính thức của dự án: chia theo
chuỗi thay vì xáo trộn ngẫu nhiên nhằm \textbf{tránh rò rỉ dữ liệu} giữa các khung
hình liên tiếp của cùng một lần bay, đồng thời buộc tập kiểm thử rơi vào các tổ hợp
góc--độ cao chưa xuất hiện y hệt trong huấn luyện (đánh giá khả năng tổng quát hóa).
Tập kiểm thử \{seq6, seq7\} gồm seq6 ($-90^\circ$, 50 m) và seq7 ($0^\circ$, 80 m).

\noindent\textbf{Lưu ý khi diễn giải.} Vì seq7 chụp ngang ($pitch = 0^\circ$) nên
tập kiểm thử \emph{có} lớp Sky --- đó là lý do IoU của Sky rất cao ($\approx 0{,}99$)
ở mọi thí nghiệm. Điều quan trọng nhất là cả bốn thí nghiệm dùng \textbf{cùng một
cách chia và cùng tập kiểm thử}, nên các con số \textbf{so sánh trực tiếp được với
nhau}. (Một số tài liệu khảo sát ban đầu của dự án từng dùng tập kiểm thử seq9 ---
thẳng xuống, không có Sky; do đó không nên đối chiếu trực tiếp các con số giữa hai
cách chia.)

\section{Kiến trúc và cấu hình huấn luyện}

Để cô lập ảnh hưởng của hàm mất mát, mọi thí nghiệm cố định kiến trúc và siêu tham
số huấn luyện như Bảng~\ref{tab:config}.

\begin{table}[H]
\centering
\caption{Cấu hình huấn luyện dùng chung.}
\label{tab:config}
\begin{tabular}{l l}
\toprule
\textbf{Thành phần} & \textbf{Giá trị} \\
\midrule
Kiến trúc            & DeepLabV3+, SegFormer, UNet, UNet++ \\
Encoder              & ResNet50 (CNN) / MiT-B4 (Transformer), khởi tạo ImageNet \\
Kích thước ảnh       & $1024 \times 1024$ \\
Batch size           & 4 \\
Optimizer            & AdamW (weight decay $10^{-4}$) \\
Learning rate        & $10^{-4}$ \\
Scheduler            & CosineAnnealingLR ($\eta_{min}=10^{-6}$) \\
Mixed precision      & Có (AMP) \\
Số lớp               & 11 (mIoU tính trên 10 lớp, bỏ \textit{Empty}) \\
Tiêu chí chọn mô hình& mIoU cao nhất trên tập kiểm định \\
\bottomrule
\end{tabular}
\end{table}

\subsection{Tăng cường dữ liệu}
Trên tập huấn luyện áp dụng các phép tăng cường (thư viện \texttt{albumentations}):
lật ngang/dọc, xoay $\pm30^\circ$, thay đổi độ sáng--tương phản, biến đổi màu
(hue/saturation), thêm nhiễu Gauss, và chuẩn hóa theo thống kê ImageNet. Tập kiểm
định và kiểm thử \textbf{chỉ} resize và chuẩn hóa (không tăng cường) để đánh giá khách
quan. Tăng cường giúp mô hình bền hơn với biến thiên góc chụp/ánh sáng và giảm
quá khớp --- đặc biệt hữu ích cho các lớp hiếm.

\subsection{Quy trình và mã nguồn}
Toàn bộ quy trình được tổ chức theo bốn notebook: \textbf{NB01} chuẩn bị gói thư viện
và trọng số pretrained offline; \textbf{NB02} huấn luyện và lưu top-$k$ checkpoint
theo mIoU kiểm định; \textbf{NB03} đánh giá trên tập kiểm thử (mIoU, PA, mF1,
per-class IoU, ma trận nhầm lẫn); \textbf{NB04} tổng hợp bảng và hình cho báo cáo. Mã
nguồn dùng chung được module hóa trong thư mục \texttt{src/} (xem
Phụ lục~\ref{ap:reproduce}). Bốn thí nghiệm của báo cáo nằm trong
\texttt{notebooks/executed/}. Luồng dữ liệu giữa bốn notebook minh họa ở
Hình~\ref{fig:pipeline}.

\begin{figure}[H]
\centering
\resizebox{\textwidth}{!}{%
\begin{tikzpicture}[
  >=stealth, font=\footnotesize, node distance=0.5cm,
  nb/.style={draw, rounded corners=3pt, minimum width=2.9cm, minimum height=1.7cm,
             align=center, fill=blue!10},
  art/.style={draw, dashed, rounded corners=2pt, align=center, fill=gray!8,
              font=\scriptsize, minimum height=0.8cm},
  data/.style={draw, cylinder, shape aspect=0.3, aspect=0.3, rotate=90,
               minimum height=2.2cm, minimum width=1.3cm, fill=green!12, align=center,
               font=\scriptsize}]
  % 4 notebook
  \node[nb] (n1) at (0,0)   {\textbf{NB01}\\Chuẩn bị\\assets offline\\{\scriptsize(Net ON, GPU OFF)}};
  \node[nb] (n2) at (4,0)   {\textbf{NB02}\\Huấn luyện\\{\scriptsize(Net OFF, GPU ON)}};
  \node[nb] (n3) at (8,0)   {\textbf{NB03}\\Đánh giá\\{\scriptsize(Net ON, GPU ON)}};
  \node[nb] (n4) at (12,0)  {\textbf{NB04}\\Báo cáo\\{\scriptsize(CPU)}};
  % artefacts giữa các bước
  \node[art, below=0.7cm of n1] (a1) {offline\_assets/\\(gói + weights)};
  \node[art, below=0.7cm of n2] (a2) {checkpoints/\\*.pth};
  \node[art, below=0.7cm of n3] (a3) {results\_all.csv};
  \node[art, below=0.7cm of n4] (a4) {bảng + hình};
  % dataset
  \node[data] (ds) at (6,2.4) {Dataset rừng\\seq1--9};
  % luồng chính
  \draw[->, thick] (n1)--(n2); \draw[->, thick] (n2)--(n3); \draw[->, thick] (n3)--(n4);
  % artefact xuống/lên
  \draw[->] (n1)--(a1); \draw[->] (n2)--(a2); \draw[->] (n3)--(a3); \draw[->] (n4)--(a4);
  \draw[->, gray] (a1) -| (n2); \draw[->, gray] (a2) -| (n3); \draw[->, gray] (a3) -| (n4);
  % dataset -> NB02, NB03
  \draw[->, green!45!black] (ds) -| (n2);
  \draw[->, green!45!black] (ds) -| (n3);
\end{tikzpicture}}
\caption{Pipeline bốn notebook: NB01 chuẩn bị gói và trọng số pretrained offline; NB02
huấn luyện (lưu checkpoint); NB03 đánh giá trên tập kiểm thử (xuất CSV); NB04 tổng hợp
bảng và hình. Dataset rừng cấp dữ liệu cho NB02 và NB03.}
\label{fig:pipeline}
\end{figure}

\section{Thiết kế thí nghiệm}
\label{sec:experiment_design}

Đề tài thực hiện nghiên cứu so sánh đối chiếu giữa hai chiến lược xử lý mất cân bằng lớp khác nhau (Bảng~\ref{tab:experiments}). Cơ sở lý thuyết của từng hàm mất mát đã được trình bày ở Mục~\ref{sec:loss}.

\begin{table}[H]
\centering
\caption{Hai cấu hình thí nghiệm so sánh chiến lược xử lý mất cân bằng lớp.}
\label{tab:experiments}
\small
\begin{tabular}{c l l c c}
\toprule
\textbf{TN} & \textbf{Chiến lược} & \textbf{Hàm mất mát} & \textbf{Epochs} & \textbf{Seed} \\
\midrule
Exp 2 & CE + trọng số lớp & PT.~\eqref{eq:ce}--\eqref{eq:enet_w} & 50 & 123 \\
Exp 4 & Focal+Dice+CE (thích nghi) & PT.~\eqref{eq:adaptive} & 50 & 123 \\
\bottomrule
\end{tabular}
\end{table}

\noindent Các notebook tương ứng của 4 kiến trúc (UNet, UNet++, DeepLabV3+, SegFormer) đặt trong các thư mục tương ứng dưới thư mục \texttt{sunny+overcast/}: \texttt{deeplabv3plus/}, \texttt{segformer/}, \texttt{unet/}, và \texttt{unet++/}.

\subsection{Chi tiết cấu hình hàm mất mát}
Để bảo đảm tính lặp lại và so sánh công bằng giữa các mô hình, cấu hình cụ thể của hai hàm mất mát như sau:
\begin{itemize}[leftmargin=1.4cm]
  \item \textbf{Exp 2 --- CE + trọng số lớp:} Sử dụng hàm mất mát Cross-Entropy kết hợp với trọng số ENet dựa trên PT.~\eqref{eq:enet_w}; trọng số thực tế được tính toán trên toàn bộ tập huấn luyện kết hợp (cả hai miền thời tiết) được cho trong Bảng~\ref{tab:weights}. Lớp nền \textit{Empty} được gán bỏ qua (\texttt{ignore\_index}).
  \item \textbf{Exp 4 --- Focal+Dice+CE (thích nghi):} Kết hợp Focal Loss ($\gamma=2{,}0$), Dice Loss (loại bỏ lớp \textit{Empty}), và Weighted CE. Điểm khác biệt cốt lõi là tỷ trọng của ba thành phần loss này ($\lambda_F, \lambda_D, \lambda_{CE}$) không cố định mà được học tự động thông qua tối ưu hóa đa nhiệm học thích nghi (Uncertainty Weighting) theo PT.~\eqref{eq:adaptive}. Các giá trị $\log\sigma_k$ được khởi tạo bằng $0$ (tương ứng trọng số ban đầu bằng $1$) và được tối ưu hóa đồng thời cùng với các trọng số của mô hình.
\end{itemize}

\begin{table}[H]
\centering
\caption{Trọng số lớp (ENet, $\kappa=1{,}02$) tính trên tập huấn luyện kết hợp (sunny+overcast), dùng cho Exp~2 và Exp~4.}
\label{tab:weights}
\small
\begin{tabular}{l r r}
\toprule
\textbf{Lớp} & \textbf{Tần suất (\%)} & \textbf{Trọng số $w_c$} \\
\midrule
Sky                & 15{,}21 & 0{,}2260 \\
Deciduous\_trees   & 16{,}36 & 0{,}2128 \\
Coniferous\_trees  & 28{,}33 & 0{,}1355 \\
Fallen\_trees      & 0{,}57  & 1{,}4142 \\
Dirt\_ground       & 1{,}99  & 0{,}9162 \\
Ground\_vegetation & 35{,}41 & 0{,}1129 \\
Rocks              & 1{,}88  & 0{,}9421 \\
Building           & 0{,}19  & 1{,}6558 \\
Fence              & 0{,}01  & 1{,}8006 \\
Car                & 0{,}05  & 1{,}7720 \\
\bottomrule
\end{tabular}
\end{table}
{\footnotesize Trọng số được chuẩn hóa để tổng bằng số lớp thực tế; các lớp cực hiếm (Fence, Car, Building) nhận trọng số lớn nhất nhằm phạt mạnh các lỗi nhận diện trên những lớp này.}

\section{Môi trường và quy trình thực nghiệm}

Quá trình huấn luyện mô hình được thực hiện trên môi trường GPU chuyên dụng (NVIDIA Kaggle / RTX 6000). Quy trình thực nghiệm bao gồm: thu thập và ghép cặp ảnh-nhãn theo chuỗi thời gian của hai miền thời tiết, tính toán trọng số ENet cho các lớp trên tập huấn luyện kết hợp, chạy huấn luyện mô hình trong 50 epoch với cơ chế lưu các checkpoint tốt nhất dựa trên mIoU của tập kiểm định, và cuối cùng là thực hiện đánh giá toàn diện mô hình trên tập kiểm thử bằng checkpoint tối ưu nhất thu được.
